\documentclass[11pt,a4paper]{article}
% preamble.tex — shared setup for human-verification documents
% Compile each document with:  tectonic <File>.tex

\usepackage[T1]{fontenc}
\usepackage[utf8]{inputenc}
\usepackage{lmodern}
\usepackage{amssymb,amsmath,amsthm}
\usepackage{mathtools}
\usepackage[dvipsnames]{xcolor}
\usepackage{listings}
\usepackage{hyperref}
\usepackage[margin=2.5cm]{geometry}
\usepackage{parskip}

% ── Colours ──────────────────────────────────────────────────────────
\definecolor{leanBg}{HTML}{F7F7F0}
\definecolor{haskellBg}{HTML}{F0F4F7}
\definecolor{keyword}{HTML}{7B3294}
\definecolor{comment}{HTML}{5A8C32}
\definecolor{noteFg}{HTML}{B35900}

% ── Lean 4 listings language ─────────────────────────────────────────
% NOTE: Lean listing blocks use ASCII approximations of unicode symbols
% (e.g. -> for arrow, o. for compose).  The literate key substitutes
% these with proper math-mode renderings in the PDF.
\lstdefinelanguage{Lean4}{
  morekeywords={def,theorem,structure,where,namespace,end,import,fun,
                let,have,show,by,exact,intro,cases,rfl,subst,obtain,
                funext,simp,rw,apply,variable,Type,Prop,instance,
                section,open,class,extends,deriving,sorry,noncomputable},
  sensitive=true,
  morecomment=[l]{--},
  morecomment=[n]{/-}{-/},
  morestring=[b]",
}

% ── Lean listing style ───────────────────────────────────────────────
\lstdefinestyle{lean}{
  language=Lean4,
  backgroundcolor=\color{leanBg},
  basicstyle=\small\ttfamily,
  keywordstyle=\color{keyword}\bfseries,
  commentstyle=\color{comment}\itshape,
  breaklines=true,
  frame=single,
  framerule=0.4pt,
  xleftmargin=4pt,
  xrightmargin=4pt,
  aboveskip=6pt,
  belowskip=4pt,
  literate=
    {->}{{$\to$}}2
    {=>}{{$\Rightarrow$}}2
    {<|}{{$\langle$}}1
    {|>}{{$\rangle$}}1
    {==>}{{$\Longrightarrow$}}3
    { o. }{{$\;\circ\;$}}3
    {forall}{{$\forall$}}1
    {:=}{{$\,:=\,$}}2
    {!=}{{$\neq$}}2
    {<=}{{$\leq$}}2
    {>=}{{$\geq$}}2,
}

% ── Haskell listing style ───────────────────────────────────────────
\lstdefinestyle{haskell}{
  language=Haskell,
  backgroundcolor=\color{haskellBg},
  basicstyle=\small\ttfamily,
  keywordstyle=\color{keyword}\bfseries,
  commentstyle=\color{comment}\itshape,
  breaklines=true,
  frame=single,
  framerule=0.4pt,
  xleftmargin=4pt,
  xrightmargin=4pt,
  aboveskip=6pt,
  belowskip=4pt,
  literate=
    {->}{{$\to$}}2
    {=>}{{$\Rightarrow$}}2,
}

% ── Theorem environments ────────────────────────────────────────────
\theoremstyle{definition}
\newtheorem{defn}{Definition}[section]
\newtheorem{thm}[defn]{Theorem}
\newtheorem{prop}[defn]{Proposition}

\theoremstyle{remark}
\newtheorem*{fidelitynote}{Fidelity note}

% ── Convenience macros ───────────────────────────────────────────────
\newcommand{\lean}[1]{\texttt{#1}}      % inline Lean identifier
\newcommand{\hask}[1]{\texttt{#1}}      % inline Haskell identifier
\newcommand{\Ext}{\mathrm{Ext}}
\newcommand{\Pos}{\mathrm{Pos}}
\newcommand{\Shape}{\mathrm{Shape}}


\title{Human Verification: \texttt{Containers.ZappaSzep}\\[4pt]
  \large Zappa--Sz\'{e}p Products: Cocycle Axioms $\Leftrightarrow$ Associativity}
\author{Machine-generated from \texttt{lean/ZappaSzep.lean} for review by Neil Ghani}
\date{July 2026}

\begin{document}
\maketitle

\begin{abstract}
This document presents every definition and theorem \emph{statement} from the
Lean~4 module \lean{Containers.ZappaSzep} in three parallel forms: the original
Lean, a readable Haskell rendering, and precise mathematical English.  Proofs
are omitted --- they are machine-checked by Lean.  The purpose is to let a
human verifier confirm that the formalisation captures the intended
mathematics.

The file formalises the \emph{pairwise Zappa--Sz\'{e}p criterion}: given two
monoids $L$ and $G$ with a left action $\triangleright : G \times L \to L$ and a
restriction $\triangleleft : G \times L \to G$, the binary operation
$(a,g) \cdot (b,h) = (a \cdot (g \triangleright b),\; (g \triangleleft b) \cdot h)$
on $L \times G$ is associative \emph{if and only if} four cocycle axioms
\lean{ZS1}--\lean{ZS4} hold (given unit compatibility).
\end{abstract}

\tableofcontents
\bigskip

%===================================================================
\section{Overview}
%===================================================================

The Zappa--Sz\'{e}p product is a non-abelian generalisation of the semidirect
product.  Where a semidirect product requires one group to act on the other, the
Zappa--Sz\'{e}p product requires a \emph{matched pair}: two \emph{mutual} actions
satisfying four compatibility axioms.

The main theorem establishes the equivalence:
$\lean{ZS1} \wedge \lean{ZS2} \wedge \lean{ZS3} \wedge \lean{ZS4}
\Leftrightarrow \mathrm{Assoc}$.
The forward direction ($\Rightarrow$) uses no unit laws; the converse
($\Leftarrow$) recovers each axiom by specialising associativity at unit elements.

%===================================================================
\section{Algebraic data}
%===================================================================

\subsection{Structure: \lean{MatchedPair}}

\begin{defn}[Matched pair of monoids]
A \textbf{matched pair} consists of:
\begin{itemize}
  \item two monoids $(L, \cdot_L, 1_L)$ and $(G, \cdot_G, 1_G)$,
  \item a left action $\triangleright : G \times L \to L$
        (``$G$ acts on $L$ from the left''), and
  \item a restriction $\triangleleft : G \times L \to G$
        (``$L$ acts on $G$ from the right''),
\end{itemize}
together with four \textbf{unit compatibility conditions}:
\begin{align*}
  \text{(U1a)}\quad & 1_G \triangleright a = a, &
  \text{(U1b)}\quad & 1_G \triangleleft a = 1_G, \\
  \text{(U2a)}\quad & g \triangleright 1_L = 1_L, &
  \text{(U2b)}\quad & g \triangleleft 1_L = g.
\end{align*}
No cocycle axiom is assumed: \lean{ZS1}--\lean{ZS4} are the \emph{content} of the
main theorem and are stated separately.
\end{defn}

\paragraph{Lean.}
\begin{lstlisting}[style=lean]
structure MatchedPair where
  L : Type
  G : Type
  mulL : L -> L -> L
  oneL : L
  mulG : G -> G -> G
  oneG : G
  mulL_assoc : forall a b c, mulL (mulL a b) c = mulL a (mulL b c)
  oneL_mul   : forall a, mulL oneL a = a
  mul_oneL   : forall a, mulL a oneL = a
  mulG_assoc : forall g h k, mulG (mulG g h) k = mulG g (mulG h k)
  oneG_mul   : forall g, mulG oneG g = g
  mul_oneG   : forall g, mulG g oneG = g
  lact       : G -> L -> L   -- notation lact g a = g |> a
  ract       : G -> L -> G   -- notation ract g a = g <| a
  lact_oneG  : forall a, lact oneG a = a
  ract_oneG  : forall a, ract oneG a = oneG
  lact_oneL  : forall g, lact g oneL = oneL
  ract_oneL  : forall g, ract g oneL = g
\end{lstlisting}

\paragraph{Haskell.}
\begin{lstlisting}[style=haskell]
data MatchedPair l g = MatchedPair
  { mulL    :: l -> l -> l
  , oneL    :: l
  , mulG    :: g -> g -> g
  , oneG    :: g
  , lact    :: g -> l -> l   -- g |> a
  , ract    :: g -> l -> g   -- g <| a
  -- plus proof obligations for associativity, units, and U1a/b, U2a/b
  }
\end{lstlisting}

\begin{fidelitynote}
In Lean, the monoid laws and unit-compatibility conditions are bundled as
\emph{fields} of the structure (proofs are data at \lean{Prop} universe).
The Haskell rendering cannot enforce them in the type; they become documentation
or runtime-checked invariants.
\end{fidelitynote}


%===================================================================
\section{The Zappa--Sz\'{e}p product}
%===================================================================

\subsection{Definition: \lean{MatchedPair.mul}}

\begin{defn}[Zappa--Sz\'{e}p product operation]
Given a matched pair $P$, define a binary operation on $L \times G$ by:
\[
  (a, g) \cdot (b, h)
  \;=\;
  \bigl(a \cdot_L (g \triangleright b),\;\; (g \triangleleft b) \cdot_G h\bigr).
\]
\end{defn}

\paragraph{Lean.}
\begin{lstlisting}[style=lean]
def MatchedPair.mul (x y : P.L * P.G) : P.L * P.G :=
  (P.mulL x.1 (P.lact x.2 y.1), P.mulG (P.ract x.2 y.1) y.2)
\end{lstlisting}

\paragraph{Haskell.}
\begin{lstlisting}[style=haskell]
zsMul :: MatchedPair l g -> (l, g) -> (l, g) -> (l, g)
zsMul p (a, g) (b, h) =
  (mulL p a (lact p g b), mulG p (ract p g b) h)
\end{lstlisting}


%===================================================================
\section{The four cocycle axioms}
%===================================================================

\subsection{Definition: \lean{MatchedPair.ZS1}}

\begin{defn}[ZS1: left action over product]
\[
  g \triangleright (a \cdot_L b)
  \;=\;
  (g \triangleright a) \cdot_L ((g \triangleleft a) \triangleright b).
\]
The left action distributes over the product in $L$, twisted by the restriction.
\end{defn}

\paragraph{Lean.}
\begin{lstlisting}[style=lean]
def MatchedPair.ZS1 : Prop :=
  forall (g : P.G) (a b : P.L),
    P.lact g (P.mulL a b) = P.mulL (P.lact g a) (P.lact (P.ract g a) b)
\end{lstlisting}

\paragraph{Haskell.}
\begin{lstlisting}[style=haskell]
zs1 :: MatchedPair l g -> Bool  -- for all g a b:
-- lact g (mulL a b) == mulL (lact g a) (lact (ract g a) b)
\end{lstlisting}


\subsection{Definition: \lean{MatchedPair.ZS2}}

\begin{defn}[ZS2: left action of $G$]
\[
  (g \cdot_G h) \triangleright a
  \;=\;
  g \triangleright (h \triangleright a).
\]
The left action is an action of $G$ (i.e.\ a monoid homomorphism $G \to \mathrm{End}(L)$).
\end{defn}

\paragraph{Lean.}
\begin{lstlisting}[style=lean]
def MatchedPair.ZS2 : Prop :=
  forall (g h : P.G) (a : P.L),
    P.lact (P.mulG g h) a = P.lact g (P.lact h a)
\end{lstlisting}

\paragraph{Haskell.}
\begin{lstlisting}[style=haskell]
-- zs2: lact (mulG g h) a == lact g (lact h a)
\end{lstlisting}


\subsection{Definition: \lean{MatchedPair.ZS3}}

\begin{defn}[ZS3: restriction over product in $G$]
\[
  (g \cdot_G h) \triangleleft a
  \;=\;
  (g \triangleleft (h \triangleright a)) \cdot_G (h \triangleleft a).
\]
\end{defn}

\paragraph{Lean.}
\begin{lstlisting}[style=lean]
def MatchedPair.ZS3 : Prop :=
  forall (g h : P.G) (a : P.L),
    P.ract (P.mulG g h) a = P.mulG (P.ract g (P.lact h a)) (P.ract h a)
\end{lstlisting}

\paragraph{Haskell.}
\begin{lstlisting}[style=haskell]
-- zs3: ract (mulG g h) a == mulG (ract g (lact h a)) (ract h a)
\end{lstlisting}


\subsection{Definition: \lean{MatchedPair.ZS4}}

\begin{defn}[ZS4: right action of $L$]
\[
  g \triangleleft (a \cdot_L b)
  \;=\;
  (g \triangleleft a) \triangleleft b.
\]
The restriction is a right action of $L$ on $G$.
\end{defn}

\paragraph{Lean.}
\begin{lstlisting}[style=lean]
def MatchedPair.ZS4 : Prop :=
  forall (g : P.G) (a b : P.L),
    P.ract g (P.mulL a b) = P.ract (P.ract g a) b
\end{lstlisting}

\paragraph{Haskell.}
\begin{lstlisting}[style=haskell]
-- zs4: ract g (mulL a b) == ract (ract g a) b
\end{lstlisting}


\subsection{Definition: \lean{MatchedPair.Assoc}}

\begin{defn}[Associativity of the Zappa--Sz\'{e}p product]
The operation $\lean{mul}$ on $L \times G$ is associative:
\[
  \forall\, x, y, z : L \times G,\quad
  (x \cdot y) \cdot z = x \cdot (y \cdot z).
\]
\end{defn}

\paragraph{Lean.}
\begin{lstlisting}[style=lean]
def MatchedPair.Assoc : Prop :=
  forall x y z : P.L * P.G, P.mul (P.mul x y) z = P.mul x (P.mul y z)
\end{lstlisting}

\paragraph{Haskell.}
\begin{lstlisting}[style=haskell]
assocProp :: MatchedPair l g -> Bool  -- for all x y z:
-- zsMul p (zsMul p x y) z == zsMul p x (zsMul p y z)
\end{lstlisting}


%===================================================================
\section{Main theorems}
%===================================================================

\subsection{Theorem: \lean{MatchedPair.assoc\_of\_zs}}

\begin{thm}[Forward direction: ZS1--ZS4 imply associativity]
Suppose $P$ is a matched pair and the four cocycle axioms hold.
Then the Zappa--Sz\'{e}p product is associative.

Proof sketch: the $L$-coordinate follows from ZS1 (re-associating the action over
a product) and ZS2 (the action law); the $G$-coordinate from ZS3 and ZS4.
No unit law is used.
\end{thm}

\paragraph{Lean.}
\begin{lstlisting}[style=lean]
theorem MatchedPair.assoc_of_zs
    (h1 : P.ZS1) (h2 : P.ZS2) (h3 : P.ZS3) (h4 : P.ZS4) :
    P.Assoc
\end{lstlisting}

\paragraph{Haskell.}
\begin{lstlisting}[style=haskell]
assocOfZS :: ZS1 p -> ZS2 p -> ZS3 p -> ZS4 p -> Assoc p
-- Proof: equational, L-coord from ZS1+ZS2, G-coord from ZS3+ZS4.
\end{lstlisting}


\subsection{Theorem: \lean{MatchedPair.zs\_of\_assoc}}

\begin{thm}[Converse direction: associativity implies ZS1--ZS4]
If the Zappa--Sz\'{e}p product is associative, then all four cocycle axioms hold.

Proof sketch: each axiom is recovered by specialising associativity at unit
elements.  ZS1 and ZS2 come from the $L$-coordinate of specific instances;
ZS3 and ZS4 from the $G$-coordinate.  The unit-compatibility fields of
\lean{MatchedPair} are exactly what is needed to collapse the instances.
\end{thm}

\paragraph{Lean.}
\begin{lstlisting}[style=lean]
theorem MatchedPair.zs_of_assoc (H : P.Assoc) :
    P.ZS1 /\ P.ZS2 /\ P.ZS3 /\ P.ZS4
\end{lstlisting}

\paragraph{Haskell.}
\begin{lstlisting}[style=haskell]
zsOfAssoc :: Assoc p -> (ZS1 p, ZS2 p, ZS3 p, ZS4 p)
-- Proof: specialise associativity at unit elements.
\end{lstlisting}


\subsection{Theorem: \lean{MatchedPair.zs\_iff\_assoc}}

\begin{thm}[Main theorem: equivalence]
For any matched pair $P$:
\[
  (\lean{ZS1} \wedge \lean{ZS2} \wedge \lean{ZS3} \wedge \lean{ZS4})
  \;\Longleftrightarrow\;
  \mathrm{Assoc}.
\]
\end{thm}

\paragraph{Lean.}
\begin{lstlisting}[style=lean]
theorem MatchedPair.zs_iff_assoc :
    (P.ZS1 /\ P.ZS2 /\ P.ZS3 /\ P.ZS4) <-> P.Assoc :=
  <|fun <|h1, h2, h3, h4|> => P.assoc_of_zs h1 h2 h3 h4, P.zs_of_assoc|>
\end{lstlisting}

\paragraph{Haskell.}
\begin{lstlisting}[style=haskell]
zsIffAssoc :: (ZS1 p, ZS2 p, ZS3 p, ZS4 p) <-> Assoc p
zsIffAssoc = (uncurry4 assocOfZS, zsOfAssoc)
\end{lstlisting}


%===================================================================
\section{Sanity-check example}
%===================================================================

\subsection{Definition: \lean{natDirect}}

\begin{defn}[Direct product of $(\mathbb{N},+,0)$ with itself]
The \textbf{trivial matched pair}: $L = G = \mathbb{N}$, multiplication $+$, unit
$0$, left action $g \triangleright a = a$ (trivial), restriction
$g \triangleleft a = g$ (identity).

This is the direct product: $(a,g) \cdot (b,h) = (a + b, g + h)$.
The four cocycle axioms hold trivially, and \lean{zs\_iff\_assoc} recovers
associativity of $(\mathbb{N}, +)$ --- a sanity check that \lean{MatchedPair} is
inhabited.
\end{defn}

\paragraph{Lean.}
\begin{lstlisting}[style=lean]
def natDirect : MatchedPair where
  L := Nat;  G := Nat
  mulL := Nat.add;  oneL := 0
  mulG := Nat.add;  oneG := 0
  mulL_assoc := fun a b c => Nat.add_assoc a b c
  oneL_mul   := fun a => Nat.zero_add a
  mul_oneL   := fun a => Nat.add_zero a
  mulG_assoc := fun a b c => Nat.add_assoc a b c
  oneG_mul   := fun a => Nat.zero_add a
  mul_oneG   := fun a => Nat.add_zero a
  lact := fun _ a => a
  ract := fun g _ => g
  lact_oneG := fun _ => rfl;  ract_oneG := fun _ => rfl
  lact_oneL := fun _ => rfl;  ract_oneL := fun _ => rfl
\end{lstlisting}

\paragraph{Haskell.}
\begin{lstlisting}[style=haskell]
natDirect :: MatchedPair Int Int
natDirect = MatchedPair
  { mulL = (+), oneL = 0
  , mulG = (+), oneG = 0
  , lact = \_  a -> a   -- trivial action
  , ract = \g  _ -> g   -- identity restriction
  }
\end{lstlisting}

\end{document}
